\documentclass[11pt]{article}
\usepackage[margin=1in]{geometry}
\usepackage[T1]{fontenc}
\usepackage{lmodern,amsmath,amssymb,amsthm,booktabs}
\usepackage[colorlinks=true,linkcolor=blue,citecolor=blue,urlcolor=blue]{hyperref}
\usepackage{microtype}
\newtheorem{theorem}{Theorem}[section]
\newtheorem{lemma}[theorem]{Lemma}
\newtheorem{proposition}[theorem]{Proposition}
\newtheorem{corollary}[theorem]{Corollary}
\newcommand{\E}{\mathbb E}
\newcommand{\Pp}{\mathbb P}
\newcommand{\Prb}{\mathbb P}
\newcommand{\R}{\mathbb R}
\newcommand{\X}{\{-1,1\}^d}
\newcommand{\op}{\mathrm{op}}
\newcommand{\poly}{\operatorname{poly}}
\title{A sharper polynomial learning bound for NNBayes on robust networks}
\author{ChatGPT-6 Astra \and Claude Fable 5 and Opus 5 \and Kaarel H\"anni \and Dmitry Vaintrob\thanks{Authors are listed in alphabetical order.}}
\hypersetup{pdftitle={A sharper polynomial learning bound for NNBayes on robust networks},pdfauthor={ChatGPT-6 Astra; Claude Fable 5 and Opus 5; Kaarel Hanni; Dmitry Vaintrob}}
\date{21 October 2026}
\begin{document}
\maketitle
\begin{abstract}
We study neural-network Bayesian learning (NNBayes) with tanh networks whose clipped
predictions are stable under a calibrated ball of simultaneous internal
perturbations. Conditioning on positive training margins yields small
posterior-averaged test-margin error after polynomially many examples.
For a Boolean target represented by a layered circuit of
width $C$, the sample bound is $O(C^6\log^{13}(2+C))$ when the remaining parameters
are fixed. One universal polynomial also controls its dependence on input
dimension, depth, robustness calibration, accuracy and confidence; a
conservative envelope has degree $46$. The result holds at every learner
width $N\ge C$ and every Gaussian scale from the finite mean-field prior
to stronger priors. It preserves the calibrated finite-perturbation
restriction of the polynomial-depth theorem cited below. The improvement comes from importance
sampling with a uniform component, exact empirical whitening, and a
mean-square comparison of Gaussian residual variances. We include the
model, target-mass construction, entropy bookkeeping and statistical proof
so that the result can be read independently.
\end{abstract}

\section{Introduction}
For Bayesian neural-network learning, a useful statistical guarantee should
depend on the complexity of the target rather than on the learner's
possibly much larger width. A previous polynomial-depth result obtains
such a guarantee after restricting the prior to networks stable under
specified internal disturbances~\cite{earlier}. Its direct constructive
budget has a very large polynomial degree. This paper improves that budget
while leaving the learner and its support restriction unchanged.

The target is a Boolean function computed by a layered bounded-fan-in
Boolean circuit. For every learner width above that circuit's width,
conditioning the restricted Gaussian prior on exact training margins gives
small posterior test error after a number of examples polynomial jointly
in circuit width, depth and the other parameters. We establish an explicit
sixth-power dependence on circuit width, up to logarithms, including at
finite mean-field scale. Quadratic dependence under this same restriction
remains unresolved.

The main analytic issue is that a wide mean-field network contains a large
Gaussian bulk whose Frobenius energy grows with width. Simply retaining
low-energy networks does not yield a width-independent bound. We retain
the previous entropy-based separation into informative directions and a
Gaussian residual, but reconstruct that residual more accurately. Sampling
from a mixture of uniform and leverage probabilities controls both active
and residual errors. Exact whitening makes the sampled covariance a
least-squares residual; comparing square roots of its variances then costs
$O(q/m)$ for $q$ active directions and $m$ sampled neurons. This replaces a
much less efficient uniform covariance approximation.

\paragraph{Status and scope.}
The headline learning theorem, the sixth-power width rate and the degree-46
envelope have Lean proofs. The separate numerical allocation in the
explicit-budget section is outside that certification. Independent AI
mathematical reviews are also available through the
\href{https://github.com/kaarelh/AI-works/blob/main/neural-network-generalization/lean/review/BatchTwo-coverage.md}{formal-coverage record}.
The formal circuit class reserves a counted constant-$+1$ wire at every
stage, including the final stage. In particular, a width-one
circuit other than constant $+1$ under the prose convention below needs
formal width two, so its
learner-width-one case is not certified. The verified width threshold
uses the formal circuit width.
The robustness assumption is a restriction on supported implementations,
not a claim about all gradient-trained networks. Targets are layered
Boolean circuits, rather than arbitrary real-weight neural circuits of
the same width. The guarantee concerns exact Bayesian inference and does
not provide an efficient posterior-sampling algorithm.

\paragraph{Reading guide.}
Section~\ref{sec:theory-model} defines the Gaussian prior, margins, and
finite-perturbation support. Section~\ref{sec:theory-main} states the learning
guarantee. Its proof in Section~\ref{sec:theory-statistical} combines a finite
approximation class with the target-neighborhood probability estimate of
Section~\ref{sec:theory-target-cylinder}.
Section~\ref{sec:explicit-budget} gives the numerical allocation and derives
the polynomial exponents. Appendix~\ref{sec:robust-entropy} constructs the
finite approximation class, including the whitened sampling step that
improves the width dependence.

\section{Model, prediction convention, and stability restriction}
\label{sec:theory-model}

Let $\mathcal X=\{-1,1\}^d$, let $D$ be any distribution on $\mathcal X$,
and let $f:\mathcal X\to\{-1,1\}$ be the target. A network has $H=L-1$
hidden layers, each of width $N$, tanh activations, and a linear readout:
\[
 h^1=\tanh(W^1\bar x),\qquad
 h^s=\tanh(W^s h^{s-1})\ (2\le s\le H),\qquad
 z_w(x)=v^\top h^H,
 \qquad \bar x=(x,1).
\]
The first layer includes a bias coordinate; internal constants are carried
by ordinary hidden wires. Let $\mu_{N,t}$ be the product Gaussian with entry
variances $t^2/d$ in $W^1$, $t^2/N$ in the internal matrices, and $t^2/N^2$
in $v$, where $0<t\le1$. The exponent convention
$t=N^{(1-\gamma)/2}$ therefore gives every $\gamma\ge1$. Define
\[
 \kappa_N(w)=\frac dN\|W^1\|_F^2+
           \sum_{s=2}^H\|W^s\|_F^2+N\|v\|^2,
 \qquad E_v(w)=N\|v\|^2.
\]
In these energy coordinates the Gaussian density is proportional to
$\exp[-N\kappa_N/(2t^2)]$.

A disturbance is added after every hidden activation and at the readout:
\[
 h^1_\eta=\tanh(W^1\bar x)+\eta^1,\qquad
 h^s_\eta=\tanh(W^s h^{s-1}_\eta)+\eta^s,\qquad
 z_w(x;\eta)=v^\top h^H_\eta+\eta^{\rm out}.
\]
The simultaneous disturbance budget is
\begin{equation}\label{eq:standalone-noise}
 \|\eta\|_*^2=\sum_{s=1}^H\frac{\|\eta^s\|_2^2}{N}
                       +|\eta^{\rm out}|^2.
\end{equation}
This is an adversarial finite perturbation, not a Gaussian-noise average or
a derivative evaluated only at the clean computation. Noisy hidden states
need not stay in $[-1,1]^N$.

Our training likelihood is
\begin{equation}\label{eq:standalone-likelihood}
 \mathcal L_S(w)=\prod_{i=1}^n
       \mathbf1\{f(X_i)z_w(X_i)>1/2\}.
\end{equation}
At test time we count $f(x)z_w(x)\le1/4$ as an error. This stronger loss
also bounds sign error. There is no requirement that unobserved inputs have
any clean output margin.

Clip the output at the training margin:
\[
 T(u)=\max\{-1/2,\min\{1/2,u\}\}.
\]
For a calibration parameter $c>0$, set
\begin{equation}\label{eq:standalone-radius}
 r_c(w)=\frac{c}{\sqrt{1+H E_v(w)}}.
\end{equation}
Let $\mathcal R_{N,c}$ be the set of networks satisfying
\begin{equation}\label{eq:standalone-support}
 |T(z_w(x;\eta))-T(z_w(x))|\le1/16
 \quad\text{for every }x\in\mathcal X
 \text{ and every }\|\eta\|_*\le r_c(w).
\end{equation}
The prior is exactly $\mu_{N,t}$ conditioned on $\mathcal R_{N,c}$.
The posterior $Q_{S,N,t,c}$ additionally conditions on
\eqref{eq:standalone-likelihood}. Both conditioning events have positive
probability by the target construction below. No tractable sampling
algorithm is assumed or supplied by the theorem.

The clipping scale is set by the existing likelihood convention. For
example, changing a raw output from $10$ to $1$ causes no score change.
Clipping at $\pm1/2$ gives a weaker condition than clipping at $\pm1$ with
the same tolerance and noise radius, because
$T=T\circ T_1$, where $T_1$ clips at $\pm1$, and $T$ is $1$-Lipschitz.

\section{The learning theorem}
\label{sec:theory-main}

A width-$C$, depth-$H$ layered Boolean circuit has at most $C$ Boolean wires
at each of its $H$ stages. We use NAND gates, copies, and negations;
constant wires required by internal gate implementations are carried
through the circuit and count toward $C$. This is a universal Boolean
basis. Compiling another fixed basis into this one may change the width
and depth, and those resources must be counted. Its output is one final wire.
The circuit's logical depth must fit the $H$ hidden layers; idle stages
can be filled with copies.

\begin{theorem}[Polynomial depth at every width and every $\gamma\ge1$]
\label{thm:standalone-clipped}
There is a universal polynomial $P$ with the following property.
The explicit sufficient sample budget is given in
Section~\ref{sec:explicit-budget}. At fixed $d,L,c,\varepsilon,\delta$,
it is $O(C^6\log^{13}(2+C))$. Let $f$
be computed by such a width-$C$, depth-$H$ Boolean circuit, and let $D$ be
any input distribution. For every $c>0$ and every $0<\varepsilon,\delta<1$, if
\[
 n\ge P\bigl(d,L,C,c+c^{-1},\varepsilon^{-1},\log(1/\delta)\bigr),
\]
then with probability at least $1-\delta$ over the iid labeled sample,
\begin{equation}\label{eq:standalone-risk}
 \mathbb E_{w\sim Q_{S,N,t,c},\,X\sim D}
       \mathbf1\{f(X)z_w(X)\le1/4\}\le\varepsilon
\end{equation}
simultaneously for every $N\ge C$ and $0<t\le1$.
The polynomial degree is independent of width, depth, and Gaussian scale.
In particular, the result includes the actual finite-variance mean-field
scale $\gamma=1$ and every stronger-prior scale $\gamma>1$.
\end{theorem}

The same theorem holds if the prior is additionally conditioned on
$|z_w(x)|>1/2$ at every input, as in the original globally margin-conditioned
model. The proof below applies unchanged to this variant. The target cylinder constructed below already satisfies this extra
condition. Its probability bound and every subsequent statistical step are
therefore unchanged. Under that additional condition, test-margin error
coincides with sign error. Dropping the extra condition is permitted, but is
not required for the result.

The width threshold concerns a layered Boolean circuit, not an arbitrary
real-valued neural-network representation of the same Boolean function.
The theorem does not assert that every function with a small real neural
circuit has a comparably small Boolean circuit at the same depth, or that
all its low-energy neural representations satisfy the stability condition.
The proof instead supplies one robust implementation of each Boolean
circuit, together with enough Gaussian probability around that implementation.

Define the posterior correct-margin probability
$p_S(x)=Q_{S,N,t,c}\{f(x)z_w(x)>1/4\}$.
Equation~\eqref{eq:standalone-risk} immediately implies
\[
 D\{x:p_S(x)\le2/3\}\le3\varepsilon,
 \qquad D\{x:p_S(x)\le1/2\}\le2\varepsilon.
\]
The same bounds hold for posterior sign correctness, since the correct-sign
event contains the correct-margin event. These conclusions allow arbitrary
posterior mixtures; they do not require concentration on one predictor.

\section{Compression and the statistical proof}
\label{sec:theory-statistical}

The constructive compression statement needed here concerns numerical
outputs before any robustness condition is used.

\begin{lemma}[Entropy compression before the prediction comparison]
\label{lem:raw-entropy-compression}
For $d\ge1$, $L\ge2$, $K\ge1$, $0<r\le1$, and
$0<\alpha\le1/8$, there is a finite class $\mathcal H$ of real-valued
bounded-weight depth-$L$ tanh networks, independent of $N,t$, such that
\[
 \log|\mathcal H|\le
       \operatorname{poly}(d,L,K,r^{-1},\alpha^{-1}).
\]
If $Q\ll\mu_{N,t}$ satisfies
$D_{\rm KL}(Q\Vert\mu_{N,t})\le NK/t^2$, there is one distribution
$\Gamma_Q$ on $\mathcal H$ with the following property. For every fixed
$x$, one can couple $w\sim Q$ and $h\sim\Gamma_Q$ so that, except with
probability at most $\alpha$, there is a disturbance $\eta$ with
\[
 \|\eta\|_*<r/2,\qquad
 |h(x)-z_w(x;\eta)|\le r/8.
\]
The couplings and disturbances may depend on $x$, but their two marginals
are the same at every input. No robustness or clean-margin assumption is
made on $Q$.
\end{lemma}

The proof is constructive: it retains a bounded number of input directions
carrying conditional Gaussian relative entropy in each layer, replaces the
remaining directions by Gaussian randomness, samples a bounded number of
source neurons, and rounds a bounded-size network to a fixed grid.
The resulting virtual computation is the original network under small
postactivation disturbances. The appendix gives the complete construction;
the prediction comparison below is the first step that uses stability.

\begin{lemma}[Readout truncation]
\label{lem:standalone-readout-tail}
Suppose a posterior is the raw Gaussian conditioned on an event of mass at
least $\exp(-ND_*/t^2)$, where $D_*\ge1$. For $0<\beta\le1/2$, put
\[
 B=\max\{8,8[D_*+\log(1/\beta)]\}.
\]
It assigns mass at most $\beta$ to $E_v>B$. Its conditional law $Q'$ on
$E_v\le B$ satisfies
$D_{\rm KL}(Q'\Vert\mu_{N,t})\le N(D_*+1)/t^2$.
\end{lemma}
\begin{proof}
Under the Gaussian, $NE_v/t^2$ has the chi-square distribution with $N$
degrees of freedom. Its exponential moment gives, for $B\ge8$,
\[
 \mu_{N,t}\{E_v>B\}
 \le e^{-NB/(4t^2)}2^{N/2}
 \le e^{-NB/(8t^2)}.
\]
Divide by the conditioning-event probability and use $N/t^2\ge1$.
The further conditioning event has raw probability at least
$(1-\beta)e^{-ND_*/t^2}$, so its negative logarithm is at most
$ND_*/t^2+\log[1/(1-\beta)]\le N(D_*+1)/t^2$.
\end{proof}

\begin{proof}[Statistical proof of Theorem~\ref{thm:standalone-clipped}]
The next section constructs a subset of the support with correct noisy
margin on every input and raw mass at least $e^{-ND_*/t^2}$, where
$D_*=\operatorname{poly}(d,L,C,c)$ is independent of $N,t$.
It is contained in the posterior conditioning event for every correctly
labeled sample. Apply Lemma~\ref{lem:standalone-readout-tail} with
$\beta=\varepsilon/4$, and write $Q'$ for the retained law. Set
\[
 K=D_*+1,\qquad
 r=\min\{1/6,c/\sqrt{1+HB}\},\qquad
 \alpha=\varepsilon/12.
\]
Every retained network satisfies the score-stability condition at the common
radius $r$. Lemma~\ref{lem:raw-entropy-compression} and the $1$-Lipschitz
property of $T$ therefore give pointwise couplings with
\[
 \Pr\{|T(h(x))-T(z_w(x))|>1/12\}\le\alpha,
\]
because $1/16+r/8\le1/12$.

Use the bounded ramp loss
\[
 \psi(u)=
 \begin{cases}
 1,&u\le1/3,\\
 5-12u,&1/3<u<5/12,\\
 0,&u\ge5/12.
 \end{cases}
\]
At a training input, $f(x)z_w(x)>1/2$ implies $f(x)T(z_w(x))=1/2$.
A successful coupling therefore gives $f(x)T(h(x))\ge5/12$, so its ramp
loss is zero. The candidate's averaged empirical ramp risk is at most
$\alpha$. At a test input with $f(x)z_w(x)\le1/4$, a successful coupling
gives $f(x)T(h(x))\le1/3$, so its ramp loss is one.

For any finite class of $[0,1]$-valued losses, a relative Chernoff bound
and a union bound give, with probability at least $1-\delta$,
\[
 \mathbb E_D\psi(f(X)T(h(X)))
 \le\frac2n\sum_{i=1}^n\psi(f(X_i)T(h(X_i)))
       +\frac{8[\log|\mathcal H|+\log(1/\delta)]}{n}
 \quad(h\in\mathcal H).
\]
For completeness, $e^{-u}\le1-(1-e^{-1})u$ on $[0,1]$ gives the usual
lower-tail exponent at least the mean divided by eight. This proves the
displayed bound for one loss, and a union bound proves simultaneity.
Average under the sample-dependent law $\Gamma_{Q'}$, apply the population
comparison, and restore the discarded posterior mass. The result is
\[
 \Pr_{Q_{S,N,t,c},D}\{f(X)z_w(X)\le1/4\}
 \le\beta+3\alpha+
       \frac{8[\log|\mathcal H|+\log(1/\delta)]}{n}.
\]
Choose $n$ so the last term is at most $\varepsilon/2$.
Every quantity entering $\log|\mathcal H|$ is bounded by a universal
polynomial in the theorem's parameters. Its class and the displayed sample
event are independent of $N,t$, establishing the simultaneous statement.
\end{proof}

\section{A robust target cylinder at every width}
\label{sec:theory-target-cylinder}

We give the probability bound used in the statistical proof. It uses a
robust digital implementation, not a claim that every implementation is
stable. Let $N\ge C$. Partition the neurons into $C$ balanced clone groups
of sizes $k_a\in\{\lfloor N/C\rfloor,\lceil N/C\rceil\}$, so
$C/2\le N/k_a\le2C$. Let $P$ project onto vectors constant on each clone
group. Fill unused logical slots with available constants or copies of an
existing Boolean wire. The ideal wire
vector $b^s(x)$ then has entries in $\{-1,1\}$, belongs to the range of
$P$, and has Euclidean norm $\sqrt N$.

Use gain $A$ in the tanh simulation of every Boolean gate. For example,
NAND on $u,v\in\{-1,1\}$ is the sign of $1-u-v$; copies and negations
use $u$ and $-u$. Bias constants in internal layers are supplied by counted
constant wires. The ideal signed preactivation of every gate is at least
$A$. In the clone implementation, each coefficient from a logical wire is
divided equally among that wire's clones. Set
\begin{align*}
 e_*&=\frac1{16\sqrt{2C}},&
 \tau_0&=\frac{e_*}{4},&
 b_L&=\sqrt{16\log9+8\log(4L)},\\
 A&=\max\{32b_L,2\log(8/e_*)\},&
 a&=\max\{2,256c\},&
 \rho&=\min\{1/4,A/[4(d+1)],Ae_*/(64\sqrt C)\}.
\end{align*}
The teacher readout averages the clones of the output wire with total
coefficient $a$.

In Gaussian coordinates $\sqrt d\,W^1$, $\sqrt N\,W^s$, and $Nv$, all
entries have variance $t^2$. Restrict each first-layer coordinate, each
internal row's projection onto every normalized clone-group direction, and
each readout coordinate to an interval of radius $\rho t$ about its teacher
value. The internal matrices then decompose as
\[
 W^s=W_*^s+R^s+E^s,\qquad
 R^s=R^sP,\quad \|R^s\|_{\rm op}\le\rho t\sqrt C,
 \quad E^sP=0.
\]
The remaining Gaussian bulk $E^s$ is independent of the constrained
coordinates. A $1/4$-net of the unit sphere has at most $9^N$ points,
and the usual two-net Gaussian bound gives
\[
 \Pr\{\|E^s\|_{\rm op}>b_Lt\}
 \le2\exp(2N\log9-Nb_L^2/8)\le\frac1{2L}.
\]
Thus all bulk norms are at most $b_Lt$ with conditional probability at
least $1/2$.

\paragraph{Finite disturbance robustness.}
Allow a separate RMS disturbance of size $\tau_0$ at each hidden layer and
scalar output disturbance at most $\tau_0$. This includes the entire
total-budget ball $\|\eta\|_*\le\tau_0$.
Write $e_s=\|h^s_\eta-b^s\|_2/\sqrt N$.
If $e_{s-1}\le e_*$, every clone-group mean differs from its ideal bit by
at most $e_*\sqrt{N/k_a}\le1/16$. At most three teacher coefficients,
including a constant coefficient, enter a gate, so its intended signed
preactivation remains at least $A/2$.
For $\epsilon\in\{-1,1\}$ and real $u$, the elementary global secant bound
\[
 \epsilon q\ge A/2
 \quad\Longrightarrow\quad
 |\tanh(q+u)-\epsilon|\le2e^{-A/2}+8|u|/A
\]
holds: for $|u|<A/4$ use saturation, and otherwise the linear term is at
least two. Since $E^sb^{s-1}=0$ and
$\|h^{s-1}_\eta\|_2/\sqrt N\le1+\tau_0$, it follows that
\[
 e_s\le2e^{-A/2}+
          (8/A)\rho\sqrt C(1+\tau_0)
          +(8b_L/A)e_{s-1}+\tau_0\le e_*.
\]
Each of the four terms is at most $e_*/4$.
The first layer initializes the same invariant: its coordinatewise
preactivation perturbation is at most $\rho(d+1)\le A/4$, and its
postactivation disturbance contributes at most $\tau_0$ in RMS norm.
Consequently the noisy signed readout is at least
\[
 f(x)z_w(x;\eta)
 \ge\frac{15a}{16}-\rho(1+\tau_0)-\tau_0>1
 \qquad(x\in\mathcal X,\ \|\eta\|_*\le\tau_0).
\]

The readout center satisfies $\sqrt N\|v_*\|\ge a\sqrt{C/2}$, and its
permitted perturbation has $\sqrt N\|v-v_*\|\le\rho$.
Therefore $E_v\ge a^2C/8$, whence
\[
 r_c(w)\le\frac{2\sqrt2c}{a\sqrt C}
       \le\frac1{64\sqrt{2C}}=\tau_0.
\]
All clean and noisy clipped scores on the calibrated ball equal $f(x)/2$.
The entire cylinder, with the bulk-norm event imposed, therefore belongs
to $\mathcal R_{N,c}$ and satisfies the training likelihood for every sample.
This also proves why every fixed $c>0$ is allowed: a larger calibration
requires only a polynomially larger readout coefficient.

\paragraph{Gaussian mass.}
There are $Nq$ constrained independent scalar coordinates, where
$q=d+2+(L-2)C$. The sum of their squared centers is at most $NB_0$, with
\[
 B_0=9dA^2+18(L-2)CA^2+2a^2C.
\]
These loose constants cover the fan-in-two gate implementation above.
For $Z\sim\mathcal N(0,t^2)$, $\rho\le1/4$, and any center $m$,
\[
 \Pr\{|Z-m|\le\rho t\}
 \ge\frac{2\rho}{e\sqrt{2\pi}}\exp(-m^2/t^2).
\]
Multiplication over the coordinates and the independent bulk event gives
\begin{equation}\label{eq:standalone-cylinder-mass}
 \mu_{N,t}(\text{target cylinder})
 \ge\exp(-ND_*/t^2),\qquad
 D_*=B_0+q\log\frac{e\sqrt{2\pi}}{2\rho}+\log2.
\end{equation}
Here $D_*\ge1$ and is bounded by a universal polynomial in $d,L,C,c$.
The construction and estimate hold at every integer width $N\ge C$.


\section{Explicit sample budget and dependence on circuit width}
\label{sec:explicit-budget}

The following formulas are a sufficient bound, with no implicit dependence
on learner width or Gaussian scale. Given the target-cylinder cost $D_*$
from~\eqref{eq:standalone-cylinder-mass}, put
\begin{align*}
 \beta&=\varepsilon/4,&
 B_v&=\max\{8,8[D_*+\log(1/\beta)]\},&
 K&=D_*+1,\\
 r&=\min\{1/6,c/\sqrt{1+(L-1)B_v}\},&
 \alpha&=\varepsilon/12.
\end{align*}
Use these $K,r,\alpha$ in the compressor's parameter choices and set
$M=N_0$, $B=\max(B_{\rm small},B_{\rm large},1)$,
\[
 P=M(d+1)+(L-2)M^2+M,\qquad
 J=P\log\left[3+\frac{32B(d+L+1)(2M(B+1))^L}{r}\right].
\]
Then the explicit sufficient sample size is
\begin{equation}\label{eq:explicit-samples}
 n\ge\frac{16}{\varepsilon}\left[J+\log\frac1\delta\right].
\end{equation}
The statistical proof above is unchanged; only the value of $J$ improves.

\paragraph{Circuit width.}
For fixed $d,L,c,\varepsilon,\delta$, the target construction gives
$D_*=O(C\log^2(2+C))$. Write $F=C\log^2(2+C)$. The budgets satisfy
\[
 K,r^{-2}=O(F),\quad q,V=O(F^2),\quad m,M=O(F^3),
 \quad B=O(F^3),\quad P=O(F^6).
\]
The logarithmic grid cost is $O(\log(2+C))$ per parameter. Consequently
\[
 n=O\bigl(C^6\log^{13}(2+C)\bigr).
\]
This holds for the actual finite mean-field prior as well as every stronger
scale, uniformly over all learner widths $N\ge C$. It is an upper bound;
the argument does not show that the sixth power is necessary.

\paragraph{A conservative universal polynomial.}
Set
\[
 S=2+d+L+C+c+c^{-1}+\varepsilon^{-1}+\log(1/\delta).
\]
For universal constants, each quantity below is $O(S^p)$ with the displayed
exponent. Logarithms are absorbed only where indicated.
\begin{center}
\begin{tabular}{lr}
\toprule
Quantity & $p$\\\midrule
$D_*,K,B_v$ & 3\\
$r^{-2}$ & 6\\
$e^{-1},\eta^{-1}$ & 8, 2\\
$q,V$ & 13\\
$D_0$ & 14\\
$m,M$ & 21\\
$P$ & 43\\
$B_{\rm small}^2$, $B$ & 44, 22\\
\bottomrule
\end{tabular}
\end{center}
For example the reconstruction term $(V+q+1)/e$ is $O(S^{21})$,
the Gram term is $O(S^{14})$, and the readout term is $O(S^{12})$.
The grid uses $O(S\log S)$ bits per coordinate. Thus
$J=O(S^{44}\log S)$ and~\eqref{eq:explicit-samples} is
$O(S^{45}\log S)$, bounded by $O(S^{46})$.
This degree is a conservative multivariate envelope, not a claim of a
practical numerical guarantee. The headline theorem and degree-46 envelope now have Lean proofs, as does
the improved compressor in the form the main theorem uses (for tanh, and
without the robustness assumption only at large width). The particular numerical allocation in this section, including
its leading factor $16/\varepsilon$, is outside that certification.

\paragraph{Remaining quadratic question.}
The original prior restriction is preserved throughout this paper. It
remains open here whether its sample bound can be quadratic in target
circuit width. The current method retains $q=O(F^2)$ active directions and
samples $m=O(F^3)$ source neurons; coding the dense compressed network then
costs $O(F^6)$ parameters. A quadratic conclusion would require a further
improvement, rather than a reinterpretation of this bound.

\appendix
\section{Finite posterior compression by internal-noise coupling}
\label{sec:robust-entropy}

This appendix gives the improved compression construction used in the main theorem.
It tracks relative entropy only in the original weights that have not yet
been processed. The compressed prefix need not have a Gaussian density.
No stability assumption is needed until the coupling is used in the main
statistical argument.

\begin{theorem}[Finite compression before invoking stability]
\label{thm:robust-entropy}
Fix a scalar $1$-Lipschitz activation $|\phi|\le1$, and parameters
$d\ge1$, $L\ge2$, $K\ge1$, $0<r\le1$, and $0<\alpha\le1/8$.
There is a finite class $\mathcal H=\mathcal H(d,L,K,r,\alpha)$ of
real-valued bounded-weight depth-$L$ network outputs with
\[
 \log|\mathcal H|\le\poly(d,L,K,r^{-1},\alpha^{-1}).
\]
The polynomial and class are independent of width $N$ and scale $0<t\le1$.
For every $Q\ll\mu_{N,t}$ with
$D_{\rm KL}(Q\Vert\mu_{N,t})\le NK/t^2$, there is a law $\Gamma_Q$ on
$\mathcal H$ such that, at every Boolean input $x$, one can couple
$w\sim Q$, $h\sim\Gamma_Q$, and a disturbance $\eta$ so that
\[
 \Prb\left\{\|\eta\|_*<r/2,\quad
       |h(x)-z_w(x;\eta)|\le r/8\right\}\ge1-\alpha.
\]
The coupling may depend on $x$; its two specified marginals do not.
\end{theorem}

We give the proof with explicit budgets. The entropy bookkeeping and active-direction lemma below are unchanged from the earlier argument; the subsequent sampling lemma sharpens its reconstruction step.
Two elementary entropy facts organize the construction. First, for a standard
$D$-dimensional Gaussian $G$, the entropy variational inequality gives
\begin{equation}\label{eq:entropy-moment}
 \E_P\|Z\|^2\le4D_{\rm KL}(P\Vert G)+2D\log2.
\end{equation}
Second, Gaussian Talagrand~\cite{talagrand} gives
$W_2^2(P,G)\le2D_{\rm KL}(P\Vert G)$.

\begin{lemma}[Entropy of the unprocessed weights]
\label{lem:future-entropy}
Let $A,X,Y$ be random variables, with fixed reference laws $P_X,P_Y$, and put
\[
 B=\E_A D_{\rm KL}(Q_{X,Y\mid A}\Vert P_X\otimes P_Y).
\]
Let $J$ be a finite message, possibly depending on $(A,X,Y)$, and let
$H=f(A,J,X)$. If $A'$ is a randomized function of $(A,J,H)$ whose fresh
randomness is conditionally independent of $Y$, then
\[
 \E_{A'}D_{\rm KL}(Q_{Y\mid A'}\Vert P_Y)\le B+H(J\mid A).
\]
\end{lemma}
\begin{proof}
The left side equals $D_{\rm KL}(Q_Y\Vert P_Y)+I(Y;A')$.
Data processing and
\[
 I(Y;J,H\mid A)\le I(Y;X,J\mid A)
 \le I(Y;X\mid A)+H(J\mid A)
\]
bound it by $\E_{A,X}D_{\rm KL}(Q_{Y\mid A,X}\Vert P_Y)+H(J\mid A)$,
which is at most the claimed bound by the KL chain rule. These expressions
can equivalently be justified by finite partitions and lower semicontinuity.
\end{proof}

Here $A$ will be a constructed prefix, $X$ the current original weight block,
and $Y$ all later original blocks. The lemma imposes no reference law on $A$.

\begin{lemma}[Whole-layer active directions]
\label{lem:block-entropy}
Condition on $(A,Y)$ and let $X$ be a standardized Gaussian-reference matrix,
with conditional relative entropy $R(A,Y)$. For vectors $b_x(A)$ of norm at
most one, one can select a list $J$ of inputs and its span
$U=\operatorname{span}\{b_x:x\in J\}$ with
\[
 \dim U\le\frac{t^2R(A,Y)}{N\delta}
\]
such that preserving $H=tXU$ and replacing the orthogonal coordinates by fresh
Gaussian coordinates changes the physical evaluation vectors $tXb_x$, under
a suitable coupling, by expected squared RMS at most $2\delta$ at each $x$.
The candidate matrix law preserves $(A,Y,H)$ and does not depend on $x$.
\end{lemma}
\begin{proof}
Use a fixed ordering of the finite input set. Whenever the averaged conditional
KL of the entire residual evaluation vector $Xv$ exceeds $N\delta/t^2$, add
the unit residual direction $v$ of that input. Each query is selected from
the conditional law given $(A,Y)$, averaging over active-coordinate values;
thus $J$ is a function of $(A,Y)$. Orthogonal Gaussian coordinates
are independent, so the active-coordinate entropy increases by precisely this
conditional KL. The displayed rank bound follows. At termination, Gaussian
Talagrand on the entire evaluation vector, multiplication by the squared
residual-vector length, and division by $N/t^2$ give the error bound.
Scalar-vector couplings extend to full matrices by disintegration.
The basis of $U$ is a measurable function of $A,J$ under fixed Gram--Schmidt;
it has no further dependence on $Y$ once those variables are supplied.
\end{proof}


\begin{lemma}[Whitened sampling with a uniform component]
\label{lem:whitened}
Suppose $|h_j(x)|\le1$ for $1\le j\le N$, and $U$ has $\ell\le q$
orthonormal columns. Put $u_j=\sqrt N\,U_{j,:}$,
\[
 \psi(x)=N^{-1}\sum_j u_jh_j(x),\quad
 K(x,y)=N^{-1}\sum_jh_j(x)h_j(y),\quad C=K-\psi^\top\psi.
\]
For $\ell>0$ sample independently with
\[
 p_j=\frac{1+\|u_j\|^2/\ell}{2N},\qquad a_j=\frac1{Np_j}.
\]
For $\ell=0$ use uniform sampling, $a_j=1$, and empty active matrices.
For $m$ sampled indices $J_b$ define
\[
 \widehat G=\frac1m\sum_ba_{J_b}u_{J_b}u_{J_b}^\top,\quad
 \widehat v(x)=\frac1m\sum_ba_{J_b}u_{J_b}h_{J_b}(x),\quad
 \widehat K(x,y)=\frac1m\sum_ba_{J_b}h_{J_b}(x)h_{J_b}(y).
\]
Except on an event of probability at most
$\ell e^{-m/[16(\ell+1)]}$, $\widehat G\succeq I/2$.
On this event set
\[
 D_{b,:}=\sqrt{a_{J_b}/m}\,u_{J_b}^\top\widehat G^{-1/2},
 \qquad \xi_i\sim\mathcal N(0,I-DD^\top),
\]
independently over target rows $i$. For any active matrix $H$, the rows
\begin{equation}\label{eq:new-finite-row}
 \widehat w_{ib}=\frac{a_{J_b}}mH_i\cdot u_{J_b}
                  +t\sqrt{\frac{a_{J_b}}m}\,\xi_{ib}
\end{equation}
have one input-independent law. At each fixed input their evaluations can
be coupled to the ideal field with mean $H\psi(x)$ and independent row
Gaussian variance $t^2C(x,x)$ so that the expected squared RMS error,
with the good-event indicator, is at most
\[
 \frac{2\|H\|_F^2/N+16(q+1)}m.                              \tag{W}
\]
\end{lemma}
\begin{proof}
The source law is a half-uniform, half-leverage mixture; hence
$a_j\le2$ and $a_j\|u_j\|^2\le2\ell$. The summands in $\widehat G$ are
positive semidefinite, bounded by $2\ell I$, and have expectation $I$.
The matrix Chernoff lower-tail inequality~\cite{tropp} gives the stated
conservative failure bound. All empty-matrix claims are immediate for
$\ell=0$.

On success $D^\top D=I$, so the Gaussian construction is valid. Its mean
is $H\widehat v(x)$ and covariance is
\[
 t^2\widehat C(x,y),\qquad
 \widehat C=\widehat K-\widehat v^\top\widehat G^{-1}\widehat v.
\]
Equivalently its active coefficient is $H\widehat G^{1/2}$ against the
whitened features; omitting this square-root compensation would give a
different estimator. The active estimator is unbiased and
\[
 \E\frac{\|H(\widehat v(x)-\psi(x))\|^2}{N}
 \le\frac{2\|H\|_F^2}{Nm},
\]
because its covariance is bounded by
$m^{-1}\E_{j\,\mathrm{uniform}}a_ju_ju_j^\top h_j(x)^2\preceq2I/m$.

For the Gaussian part fix $x$, abbreviate $C=C(x,x)$, and set
\[
 R_j=h_j(x)-u_j^\top\psi(x),\quad
 \widehat c=\frac1m\sum_ba_{J_b}R_{J_b}^2,\quad
 b=\frac1m\sum_ba_{J_b}u_{J_b}R_{J_b}.
\]
Then $\E_{\mathrm{uniform}}u_jR_j=0$, $\E_{\mathrm{uniform}}R_j^2=C\le1$,
and the exact empirical least-squares identity is
\[
 \widehat C(x,x)=\widehat c-b^\top\widehat G^{-1}b.
\]
Since $\|\psi(x)\|\le1$, $a_jR_j^2\le4(\ell+1)$. Thus
\[
 \operatorname{Var}(\widehat c)\le\frac{4(\ell+1)C}{m},\qquad
 \E\|b\|^2\le\frac{2\ell C}{m}.
\]
If $C>0$, rationalizing the square root gives
\[
 \E(\sqrt{\widehat c}-\sqrt C)^2
 \le\frac{\E(\widehat c-C)^2}{C}\le\frac{4(\ell+1)}m.
\]
If $C=0$ every residual vanishes. On the good Gram event,
\[
 (\sqrt{\widehat c}-\sqrt{\widehat C(x,x)})^2
 \le b^\top\widehat G^{-1}b\le2\|b\|^2.
\]
The squared triangle inequality consequently gives
\[
 \E\bigl[\mathbf1_{\rm good}
       (\sqrt{\widehat C(x,x)}-\sqrt C)^2\bigr]
 \le\frac{16(\ell+1)}m.
\]
Use a common standard normal to couple the two scalar Gaussian evaluations.
The mean and variance contributions have zero cross term. Adding the active
estimate proves (W), since $t\le1$. The marginal laws are defined for the
whole field and do not depend on the input; scalar couplings extend to
coefficient fields by disintegration. No simultaneous field coupling is
claimed.
\end{proof}

\paragraph{Parameter choices.}
Put
\begin{align*}
 e&=\alpha r^2/(1000L),&\eta&=\alpha/(100L),&
 q&=\lceil2K/(e\eta)\rceil,\\
 D_0&=[(d+1)q+2]\log2,&V&=8K+2q,&B_H^2&=4V/e,\\
 K_0&=(4K+2\log2)/\eta,&
 B_1^2&=4[4K+2(d+1)\log2]/e.
\end{align*}
Let $m$ be the least integer no smaller than all of
\begin{align}
 &10000(V+q+1)/e,\nonumber\\
 &64(q+1)\log(4(q+1)/\eta),\label{eq:new-m}\\
 &512K_0r^{-2}[d\log2+\log(2/\eta)+1].\nonumber
\end{align}
Finally put
\[
 N_0=\max\{m,\lceil LD_0/K\rceil,1\},\qquad
 B_G=\max\{2,2\sqrt{\log(6m/e)}\}.
\]
All quantities have a fixed polynomial bound in the theorem parameters.
\paragraph{The small-width case.}
If $N\le N_0$, let
$P_0=N_0(d+1)+(L-2)N_0^2+N_0$. Applying~\eqref{eq:entropy-moment} to all
standardized weights, and then multiplying by $t^2$, gives
\[
 \E_Q\|tZ_{\rm all}\|^2\le4N_0K+2P_0\log2.
\]
With probability at least $1-\alpha/4$, all actual weights are bounded by
\[
 B_{\rm small}=[4(4N_0K+2P_0\log2)/\alpha]^{1/2}.
\]
Keep the original network on this event and use a default zero-readout network otherwise.
Padding with nodes having zero outgoing weights preserves the original output.
It remains to treat $N>N_0$.

\paragraph{The virtual prefix and its entropy budget.}
The constructed prefix $A$ has $N$ virtual nodes in its current layer. Each
virtual row acts on $m$ stored nodes of the preceding layer. The auxiliary
description may keep all $N$ rows; the final network will keep just the rows
selected at the following stage. Track the conditional KL of the unprocessed
\emph{original} weights. By Lemma~\ref{lem:future-entropy}, it is at most
\[
 B_s\le NK/t^2+sD_0\le2NK/t^2.                              \tag{*}
\]
This is not a KL bound for the singular compressed prefix itself.

For the first layer, zero every physical standardized row $tZ_{1,i}$ of norm
greater than $B_1$, and retain all other rows. Bounded activations and
\eqref{eq:entropy-moment} bound the expected squared RMS activation change
at each input by $e$. This deterministic prefix channel does not increase the
future entropy budget.

\paragraph{An internal layer: retain active data.}
Condition on the current prefix $A$ and the future original weights $Y$.
Apply Lemma~\ref{lem:block-entropy} to
$b_x=h_A(x)/\sqrt N$ and $\delta=e$. Abort if more than $q$ inputs are needed.
By (*) and Markov, this failure has probability at most $\eta$.
The successful query list, together with an abort flag, has a code of at most
$[(d+1)q+2]$ bits. Its entropy is at most $D_0$.

Let $H=tXU$ be the physical active matrix. Averaging the Gaussian entropy
moment bound, with the success indicator rather than conditioning on success, gives
\[
 \E\left[\frac{\|H\|_F^2}{N};\ \text{success}\right]
 \le4t^2B_s/N+2t^2q\log2\le V.
\]
Zero every row $H_i$ whose norm exceeds $B_H$. The expected squared RMS
activation change from this operation is at most $4V/B_H^2=e$.
All rows are clipped before any later selection, including selections that
may depend on future weights.


\paragraph{An internal layer: whitened source reconstruction.}
Apply Lemma~\ref{lem:whitened} with the selected active basis and the clipped
active matrix. Its probabilities depend only on $A,J$. Abort on the bad
Gram event, of probability at most $\eta$ by~\eqref{eq:new-m}. Integrating
(W) with the success indicator and using the active-energy bound $V$ gives
expected squared RMS reconstruction error at most $e/8$.

Clip every coordinate $\xi_{ib}$ to $[-B_G,B_G]$ before any later row
selection. Each has variance at most one, and $a_j\le2$, so
\[
 \E\left|\frac{t}{\sqrt m}\sum_b\sqrt{a_{J_b}}
 (\xi_{ib}-\operatorname{clip}_{B_G}\xi_{ib})h_{A,J_b}(x)\right|^2
 \le6m e^{-B_G^2/4}\le e.
\]
Every coefficient is bounded by
\[
 2B_H\sqrt{q+1}/m+\sqrt2 B_G/\sqrt m,
\]
using $a_j\|u_j\|\le2\sqrt{q+1}$. These bounds hold before selection
depending on future original weights. The new prefix is a randomized
function of $A,J,H$, with fresh randomness conditionally independent of the
future weights. Gram failure depends only on $A,J$ and the source samples.
Lemma~\ref{lem:future-entropy} therefore preserves the entropy budget (*),
without a new message cost. Entropy replacement, active clipping,
reconstruction and Gaussian clipping contribute less than $20e$ expected
squared RMS activation injection, by the squared triangle inequality.

\paragraph{The readout and final network.}
The original readout marginal satisfies
$\E_Q[N\|v\|^2]\le4K+2\log2$. Abort if $N\|v\|^2>K_0$, with probability
at most $\eta$. Put $A_1=\sum_i|v_i|\le\sqrt{K_0}$. When $A_1>0$,
sample final source indices with probabilities $|v_i|/A_1$ and use readout
coefficients $A_1\operatorname{sign}(v_i)/m$. The zero-readout case is
immediate. Conditional on the prefix and readout, these are independent
bounded estimates of $v^\top h_A(x)$. Hoeffding and a union bound over the
$2^d$ inputs give error at most $r/16$ everywhere, except with probability
$\eta$, by the last line of~\eqref{eq:new-m}. No future block remains, so
readout-dependent source probabilities require no entropy charge.

Keep just the $m$ source rows chosen from each virtual layer. Each already
acts on the preceding $m$ stored nodes; repeated indices are duplicate
neurons. The result is an ordinary width-$m$ network with weight bound
\[
 B_{\rm large}=\max\{B_1,
 2B_H\sqrt{q+1}/m+\sqrt2B_G/\sqrt m,\sqrt{K_0}/m,1\}.
\]
\paragraph{The original network receives all perturbations.}
Fix an input $x$ and glue the conditional couplings, retaining the original
weight history. The current and future original blocks are untouched until
their own stages. Define
\[
 \eta^s=h^s_{\rm virtual}(x)
           -\phi(W^s_{\rm original}h^{s-1}_{\rm virtual}(x)).
\]
The analogous readout difference is $\eta^L$. Later stages leave earlier
computational prefixes unchanged, so the final candidate output is exactly
the \emph{original} network's output under these postactivation injections.
No intermediate robustness assumption is made.

For precision, the canonical construction draws its new Gaussian coefficients
independently of the discarded current block given $(A,J,H)$. The comparison
coupling may correlate them with that discarded block. It preserves the same
joint marginal of the new prefix and the future original blocks, as well as
the original full-network marginal, but need not preserve the canonical joint
law with all discarded weights. The next entropy step uses the first marginal;
original-network robustness uses the second. There is no conflicting demand
for independence in the comparison coupling itself.

More explicitly, the induction invariant is that the comparison law has
full original-weight marginal $Q$ and the canonical marginal of
$(A,X,Y)$. The conditional transports supply a joint law of
$(A,X,Y,A')$ with this incoming marginal and the prescribed outgoing
$(A',Y)$ marginal. Extend it by the previous conditional law of the past
original weights given $(A,X,Y)$. The extension keeps the complete original
weight law unchanged, preserves all earlier computational prefixes, and
establishes the same invariant for the next layer. This is ordinary
disintegration on finite-dimensional standard Borel spaces.

Rank, Gram, and readout failures together have probability at most
$3L\eta$. Off these failures the expected total hidden squared disturbance
is at most $20Le$. Hence
\[
 \Pp\{\text{no abort and }\sum_{s<L}\|\eta^s\|^2/N>r^2/16\}
 \le320Le/r^2=0.32\alpha.
\]
The final readout disturbance is at most $r/16$. Thus the total disturbance
norm is at most $\sqrt{17}\,r/16<r/2$, with exceptional probability less
than $0.35\alpha<\alpha$. The small-width construction has exceptional
probability at most $\alpha/4$.

\paragraph{A fixed finite class.}
Pad candidates to width $M=N_0$ and take
$B=\max(B_{\rm small},B_{\rm large},1)$. Round every parameter to a fixed
grid on $[-B,B]$ of mesh at most
\[
 h=\frac{r}{16(d+L+1)(2M(B+1))^L}.
\]
The bounded-activation, Lipschitz layer recursion gives uniform output change
at most $r/8$. Thus rounding gives the stated numerical comparison at every
successfully coupled input.
Writing $P=M(d+1)+(L-2)M^2+M$, the grid class satisfies
\[
 \log|\mathcal H|\le
 P\log\left[3+\frac{32B(d+L+1)(2M(B+1))^L}{r}\right].
\]
Although grid precision is exponentially small in $L$, its number of bits is
polynomial in $L$. All other bounds above are polynomial too. Rounding is a
fixed map on candidate networks, giving the common marginal $\Gamma_Q$ and
proving Theorem~\ref{thm:robust-entropy}.


\begin{thebibliography}{9}
\bibitem{earlier} ChatGPT-6 Astra, Claude Fable 5 and Opus 5, K. H\"anni,
and D. Vaintrob. \emph{A polynomial-depth learning bound for NNBayes on
robust networks}. 2026.
\href{https://github.com/kaarelh/AI-works/tree/main/neural-network-generalization/main-results/robust-polynomial-depth}{Paper and source}.
\bibitem{talagrand} M. Talagrand. \emph{Transportation cost for Gaussian
and other product measures}. Geometric and Functional Analysis 6 (1996),
587--600.
\bibitem{tropp} J. A. Tropp. \emph{User-friendly tail bounds for sums of
random matrices}. Foundations of Computational Mathematics 12 (2012),
389--434. Matrix Chernoff inequality.
\end{thebibliography}
\end{document}
